% Generated draft: requires mathematical and source review.
\documentclass[11pt]{article}
\usepackage[margin=25mm]{geometry}
\usepackage{fontspec}
\setmainfont{Latin Modern Roman}
\usepackage{amsmath,amssymb,mathtools,unicode-math}
\setmathfont{Latin Modern Math}
\usepackage{enumitem,needspace,xurl,hyperref}
\setlength{\parindent}{0pt}
\setlength{\parskip}{5pt}
\setlength{\emergencystretch}{3em}
\newcommand{\N}{\mathbb{N}}
\newcommand{\Z}{\mathbb{Z}}
\newcommand{\Q}{\mathbb{Q}}
\newcommand{\R}{\mathbb{R}}
\newcommand{\C}{\mathbb{C}}
\newcommand{\F}{\mathbb{F}}
\newcommand{\A}{\mathbb{A}}
\newcommand{\PP}{\mathbb{P}}
\newcommand{\E}{\mathbb{E}}
\newcommand{\eps}{\varepsilon}
\newcommand{\dd}{\,\mathrm{d}}
\begin{document}
\section*{UnsolvedMath ID: 30001401}
\textbf{Source problem code:} catalogue entry
\subsection{Definitions and mathematical statement}
The specified dynamical group is
\[
G=\mathrm{SU}(2)=\{U\in M_2(\C):U^*U=I_2,\ \det U=1\},
\]
where \(U^*\) is the conjugate transpose. For a complex Hilbert space \(\mathcal H\), a pure state is represented by a ray \([\psi]\in\PP(\mathcal H)\), where \(\psi\ne0\) and vectors differing by a nonzero complex scalar represent the same ray. The supplied record does not specify \(\mathcal H\), a representation of \(G\) on \(\mathcal H\), or which rays are to be called classical or entangled.

\textbf{Question as supplied.} Does the Bell-inequality criterion for distinguishing classical from entangled pure states admit an analogue when the dynamical group is \(G=\mathrm{SU}(2)\)?

\textbf{Unresolved formulation.} The record does not identify the original Bell-inequality criterion, the permitted observables or measurement arrangements, the inequalities sought, or whether the desired analogue must be necessary, sufficient, or both. It also does not specify whether the question concerns a particular representation or a quantified family of representations of \(\mathrm{SU}(2)\). Thus the displayed group does not by itself determine a precise criterion or its quantifiers. The supplied literature assessment additionally describes an unresolved geometric question about linear dependence of classical states, but gives neither that formulation nor its connection to the stated Bell-inequality question. These missing specifications prevent a unique precise mathematical reformulation.
\subsection{Short English statement}
The supplied question asks for an analogue, for \(\mathrm{SU}(2)\) dynamics, of a Bell-inequality test distinguishing classical and entangled pure states. The state space, representation, meaning of classicality, and required test are unspecified. The record labels the problem open; this is supplied curation rather than an independently verified status.
\subsection{Sources}
\begin{itemize}
\item Supplied local UnsolvedMath registry record, permanent ID \(30001401\), title: \emph{Bell Inequality Criteria for Special Unitary Dynamics}; collection: Oberwolfach Reports Open Problems; exact source problem code: \texttt{catalogue entry}.
\item Supplied report citation: \emph{Mini-Workshop: Geometry of Quantum Entanglement (2010)}, Oberwolfach Reports, DOI \texttt{10.4171/owr/2009/56}, pp. 2993--3032. Supplied source URL: \url{https://doi.org/10.4171/owr/2009/56}. No individual proposer is supplied. The registry gives proposed year 2009, while the supplied report title gives 2010; this discrepancy is retained without resolution.
\item Supplied literature reference, labelled ``original Oberwolfach problem report'': \url{https://doi.org/10.4171/owr/2009/56}.
\item Literature checked at: \texttt{2026-08-21}. This is supplied dated curation, not independently verified. Its assessment describes an unresolved \(\mathrm{SU}(2)\) geometric question concerning linear dependence of classical states and reports that no later resolution was located. The background explicitly requests independent verification of time-sensitive status and cited claims.
\item The registry supplies status \texttt{open} and statement status \texttt{corrected\_verified}. Its verification note says that the complete statement was preserved while mathematical notation was placed in delimiters. These are source metadata and do not constitute independent mathematical review of this draft.
\end{itemize}
\end{document}
